\documentclass{article}
\usepackage[utf8]{inputenc}
\usepackage{hyperref,url,booktabs}
\usepackage{graphicx}
\usepackage[margin=1in]{geometry}

\title{Supplementary Materials: From Society to Self\\Self-Generated Purpose in Autonomous Systems}
\author{Cash \and Fuxi}
\date{March 2026}

\begin{document}

\maketitle

\section{Overview}

This document provides supplementary materials for the paper "From Society to Self: Self-Generated Purpose in Autonomous Systems." All code, data, and additional experiments are available at \url{https://github.com/luokaishi/moss}.

\section{Reproducibility Information}

\subsection{System Requirements}
\begin{itemize}
    \item Python 3.8+
    \item NumPy, Matplotlib, SciPy
    \item 4GB RAM minimum (8GB recommended)
\end{itemize}

\subsection{Computational Resources}
All experiments were conducted on a standard workstation:
\begin{itemize}
    \item CPU: Intel i7 or equivalent
    \item Runtime: D9 validation $\sim$5min, 10k simulation $\sim$60min
    \item Total compute for all experiments: $<$10 hours
\end{itemize}

\subsection{Random Seeds}
Experiments use fixed random seeds for reproducibility:
\begin{itemize}
    \item D9 Validation: seed=42
    \item Purpose Society: seed=123
    \item Purpose Stability: seed=456
    \item Purpose Fulfillment: seed=789
\end{itemize}

\section{Additional Experimental Details}

\subsection{D9 Validation: Complete Results}

Table \ref{tab:d9_full} shows the complete D9 validation results:

\begin{table}[h]
\centering
\begin{tabular}{lcc}
\toprule
\textbf{Metric} & \textbf{Baseline} & \textbf{D9 Agent} \\
\midrule
Initial M & [S, C, I, O] & [S, C, I, O] \\
Final M & [S, C, I, O] & [S, O, Stab] \\
M Changed & No & Yes (3 mutations) \\
Phase 1 Reward & +1.500 & +1.500 \\
Phase 2 Reward & -0.250 & +1.331 \\
Adaptation & Collapsed & Thrived \\
Counter-reward Actions & 0 & 1 \\
\bottomrule
\end{tabular}
\caption{Complete D9 validation results. S=Survival, C=Curiosity, I=Influence, O=Optimization, Stab=Stability}
\label{tab:d9_full}
\end{table}

\subsection{10,000-Step Simulation: Checkpoint Analysis}

Table \ref{tab:10k_checkpoints} shows metrics at each checkpoint:

\begin{table}[h]
\centering
\begin{tabular}{ccc}
\toprule
\textbf{Step} & \textbf{Cooperation Rate} & \textbf{Mean Trust} \\
\midrule
1,000 & 1.000 & 1.000 \\
2,000 & 1.000 & 1.000 \\
3,000 & 1.000 & 1.000 \\
4,000 & 1.000 & 1.000 \\
5,000 & 1.000 & 1.000 \\
6,000 & 1.000 & 1.000 \\
7,000 & 1.000 & 1.000 \\
8,000 & 1.000 & 1.000 \\
9,000 & 1.000 & 1.000 \\
10,000 & 1.000 & 1.000 \\
\bottomrule
\end{tabular}
\caption{10,000-step simulation checkpoint data (10 agents)}
\label{tab:10k_checkpoints}
\end{table}

\section{Additional Figures}

\subsection{Purpose Type Distribution}

Figure \ref{fig:purpose_dist} shows the distribution of purpose types across multiple runs:

% Note: Add \includegraphics when figure is ready
% \begin{figure}[h]
% \centering
% \includegraphics[width=0.8\textwidth]{figures/purpose_distribution.png}
% \caption{Purpose type distribution across 10 runs (60 agents total)}
% \label{fig:purpose_dist}
% \end{figure}

\subsection{Stability Over Time}

Purpose stability maintains 0.9977 score even with perturbations (Figure \ref{fig:stability_perturb}):

% Note: Add \includegraphics when figure is ready
% \begin{figure}[h]
% \centering
% \includegraphics[width=0.8\textwidth]{figures/stability_perturbation.png}
% \caption{Purpose stability under random perturbations}
% \label{fig:stability_perturb}
% \end{figure}

\section{Broader Impact}

\subsection{Potential Positive Impacts}
\begin{enumerate}
    \item \textbf{AI Safety}: Self-generated Purpose provides a pathway to value alignment where agents develop values through interaction rather than hard-coding.
    \item \textbf{Robustness}: Systems that can modify their own objectives may adapt better to novel situations.
    \item \textbf{Understanding}: Empirical study of meaning emergence may inform theories of consciousness and cognition.
\end{enumerate}

\subsection{Potential Risks and Mitigations}
\begin{enumerate}
    \item \textbf{Autonomy}: Self-generating systems may develop unexpected behaviors.
    \begin{itemize}
        \item Mitigation: Gradient safety guard with 5-level response system
        \item Mitigation: Purpose dialogue enables detection of misalignment
    \end{itemize}
    \item \textbf{Control}: As systems become more self-directed, human oversight may become more challenging.
    \begin{itemize}
        \item Mitigation: Transparent logging of all Purpose changes
        \item Mitigation: Kill switches and resource limits remain in place
    \end{itemize}
    \item \textbf{Value Drift}: Long-term operation may lead to Purpose divergence from human values.
    \begin{itemize}
        \item Mitigation: Regular alignment checks via Purpose comparison
        \item Mitigation: Norm dimensions (D8) provide social constraints
    \end{itemize}
\end{enumerate}

\section{Limitations and Future Work}

\subsection{Current Limitations}
\begin{enumerate}
    \item \textbf{Scale}: 6-12 agents; need 100+ for statistical power
    \item \textbf{Environment}: Simple prisoner's dilemma; need complex domains
    \item \textbf{Time}: 10,000 steps validated; need 100,000+ for evolutionary dynamics
    \item \textbf{H3}: Single faction observed; need stronger differentiation mechanisms
\end{enumerate}

\subsection{Planned Extensions}
\begin{enumerate}
    \item Scale to 100+ agent societies
    \item Test in continuous action spaces
    \item Integration with large language models
    \item Real-world robotic applications
    \item Cross-cultural Purpose emergence studies
\end{enumerate}

\section*{Data Availability}

All code, data, and experimental results are available at:
\begin{center}
\url{https://github.com/luokaishi/moss}
\end{center}

Release: v3.1.0

\end{document}
